\section{Introduction}

Demand planning increasingly relies on predictive models to anticipate future demand at fine operational granularity. In many empirical studies and model-selection workflows, the primary criterion for choosing among forecasting methods is numerical prediction accuracy. Metrics such as mean absolute error, root mean squared error, and weighted absolute percentage error are important because inaccurate forecasts can lead to costly overstocking, stockouts, and service failures. However, forecast accuracy alone does not describe whether the resulting plan can be executed.

In real operations, forecasts are not executed directly. A forecast is converted into a planning signal: an inventory target, a replenishment quantity, a labor plan, a production target, a transportation commitment, or an allocation decision. This forecast-to-plan conversion is mediated by business policies, service-level targets, safety-stock rules, system constraints, and planner review processes. The operational question is therefore not only whether a forecast is numerically close to realized demand, but whether the planning signal implied by that forecast is stable, feasible, and worth deploying.

A more accurate forecast may create unstable planning signals. It may require frequent model switching, large period-to-period target jumps, repeated infrastructure updates, or adjustment requests that exceed what suppliers, stores, warehouses, transportation teams, and enterprise planning systems can absorb. These burdens may be invisible in forecast-error tables. A model that looks attractive under WAPE can still be unattractive when it generates plan churn, execution violations, or additional governance work.

This creates a planning-execution gap. Forecasting capability can adapt faster than downstream execution infrastructure. Machine learning pipelines may update daily or weekly, while procurement lead times, labor scheduling, supplier coordination, production planning, and information-system workflows may change more slowly. When the forecasting layer changes faster than the execution layer can absorb, prediction accuracy and deployable planning quality can diverge.

This paper studies forecast-driven demand planning as a feasibility-constrained model selection problem. Rather than asking which forecasting model has the lowest error in isolation, the paper asks which deployable planning strategy balances forecast accuracy, inventory exposure, planning volatility, model-switching burden, and execution-capacity constraints. The objective is not to argue that stability-aware or feasibility-aware strategies always dominate. The objective is to make the tradeoff surface visible and to evaluate model selection in the same operational space where planning decisions are executed.

The empirical design uses three public-data roles. DataCo supply-chain execution data is used to show that execution risk is empirically measurable and to construct data-informed execution-risk sensitivity scenarios. Favorita retail demand data provides the main forecast-to-plan proof of concept, with demand, promotion, holiday, oil-price, transaction, and store-context signals. M5 provides a larger hierarchical retail robustness setting across item, department, category, store, and state levels. Walmart and Rossmann are reserved as optional future robustness extensions; no results from those datasets are claimed in the current draft.

The paper makes four contributions. First, it formulates forecast-driven demand planning as feasibility-constrained model selection, separating forecast generation from the downstream choice of a deployable planning strategy. Second, it defines a normalized planning utility that compares inventory exposure, planning volatility, switching burden, and execution violations relative to an accuracy-first baseline while still reporting raw operational metrics. Third, it constructs DataCo-informed execution-risk scenarios as sensitivity anchors, without claiming exact economic cost calibration between DataCo and the demand-planning datasets. Fourth, it shows empirically that execution-risk-sensitive objectives can reorder deployable strategy rankings and that the preferred deployable method depends on execution scenario, planning granularity, and demand intermittency.
